\section{Background}\label{sec:background}

\subsection{Tsetlin Machine}
A TM operates on a Boolean input $\mathbf{x}\in\{0,1\}^F$.
Its literals are the features $x_i$ and their negations $\lnot x_i$~\cite{granmo2018tsetlin}.
Each class owns a set of clauses. Half of them have positive polarity and half negative polarity.
A clause $C_j$ is the conjunction of its included literals $\mathcal{I}_j$.
In the strict TM, the class vote is
\begin{equation}
v_c(\mathbf{x}) = \sum_{j\in\mathcal{P}_c} C_j(\mathbf{x}) - \sum_{j\in\mathcal{N}_c} C_j(\mathbf{x}),
\end{equation}
where $\mathcal{P}_c$ and $\mathcal{N}_c$ are the positive and negative clauses of class $c$.
The predicted class is $\arg\max_c v_c$.
Training uses teams of Tsetlin automata, one per literal and clause, which decide whether the literal is included~\cite{granmo2018tsetlin}. A vote threshold $T$ sets the target margin of the class votes during training.
Inference uses only the included literal sets.
In the vanilla TM implementation used here~\cite{pytsetlin}, an empty clause outputs 0 at inference.

\subsection{Fuzzy-Pattern Tsetlin Machine}
FPTM replaces the strict conjunction with a graded score~\cite{hnilov2025fptm,hnilov2025code}.
Let $n_j=|\mathcal{I}_j|$ and let $m_j(\mathbf{x})$ be the number of included literals that are false for $\mathbf{x}$.
The reference implementation (function \texttt{check\_clause} in~\cite{hnilov2025code}) computes the clause score as
\begin{equation}\label{eq:fptm}
s_j(\mathbf{x}) = \max\bigl(0,\; b_j - m_j(\mathbf{x})\bigr),\qquad
b_j = \begin{cases} \min(n_j, \LF), & n_j>0,\\ \LF, & n_j=0,\end{cases}
\end{equation}
and the class vote is $v_c=\sum_{\mathcal{P}_c}s_j-\sum_{\mathcal{N}_c}s_j$.
$\LF$ is a hyperparameter that sets how many literal misses a clause absorbs.
With $\LF=1$, a non-empty clause scores 1 if and only if all its literals hold, which is the strict TM clause. An empty clause scores $\LF$, i.e., 1 for $\LF=1$.
FPTM keeps the TM feedback structure.
Relative to a vanilla TM, it introduces three modifications, all visible in the reference code~\cite{hnilov2025code}.
\begin{enumerate}
\item \emph{Literal budget $L$.} Type~Ia feedback may add literals to a clause only while the clause has at most $L$ included literals.
\item \emph{Deterministic Type~Ia feedback.} When a selected clause has a positive score, the automata of all true literals are rewarded toward inclusion, with no random sub-sampling. The penalty of false literals below the include threshold is deterministic as well.
\item \emph{Fuzzy clause voting} according to Eq.~(\ref{eq:fptm}). A clause counts as firing in feedback whenever its score is positive, not only on a full match.
\end{enumerate}
The reference trainer also differs from Granmo's multiclass TM in two implementation details.
First, it applies negative feedback to every class other than the target class, whereas the multiclass TM samples one negative class per example~\cite{granmo2018tsetlin}.
Second, Type~Ib feedback of a non-firing clause decrements $d=\mathrm{round}(F/S)$ randomly chosen automata among the plain literals and $d$ among the negated literals. The vanilla TM instead decrements each automaton with probability $1/s$, where $s$ is its specificity.
Modifications~1 and~2 and these two details affect training only.
At inference, FPTM differs from the vanilla TM only through modification~3. It is therefore the only one that the \gainp{} mapping must support.

\subsection{\gainp{}: a GC-eDRAM associative processor}\label{sec:gainp}
\gainp{} is an AP whose CAM array is an unmodified hybrid-3T GC-eDRAM~\cite{levi2026gainp}.
A silicon-proven GC-eDRAM macro in 16~nm FinFET technology forms its core~\cite{levi2026gainp}.

\paragraph{Gain cell}
Fig.~\ref{fig:gc}(a) shows the hybrid-3T gain cell, which combines NMOS and PMOS devices~\cite{levi2026gainp}.
The write port is a single NMOS transistor M$_\mathrm{W}$.
A write drives the write bitline (WBL) to the data value and asserts the write wordline (WWL), which charges or discharges the storage node (SN).
SN drives the gate of the PMOS transistor M$_\mathrm{S}$. M$_\mathrm{S}$ and the PMOS transistor M$_\mathrm{R}$ in series form the read port.
A read pre-discharges the read bitline (RBL) and drives the read wordline (RWL) low.
If SN holds 0, M$_\mathrm{S}$ conducts and the RBL charges toward $V_\mathrm{DD}$. If SN holds 1, the path is cut off and the RBL stays low~\cite{levi2026gainp}.
The RBL is sensed by a digital sense inverter, enabled by a timing signal from a replica bitline column~\cite{levi2026gainp}.
The stored charge leaks, mainly through subthreshold leakage of the write device and gate leakage of the read devices, so the cell needs refresh~\cite{levi2026gainp}.
The 16~nm design reduces leakage with a single-fin write transistor, a negative WWL bias of $-150$~mV for unselected rows, a boosted write supply and high-$V_\mathrm{T}$ read devices~\cite{levi2026gainp}.

\paragraph{CAM cell and search}
\gainp{} stores data in column-major order~\cite{levi2026gainp}.
Each bit occupies two vertically stacked gain cells. The top cell holds the bit value $\beta$ and the bottom cell holds $\overline{\beta}$ (Fig.~\ref{fig:gc}(b)).
The horizontal RWLs act as searchlines and the vertical RBLs as matchlines (ML).
In a compare, the inverted query bit drives the top RWL ($\overline{\mathrm{SL}}$) and the query bit drives the bottom RWL (SL).
A read port then conducts only when the query bit differs from the stored bit.
The read current is additive. A matchline therefore charges when at least one cell of the column mismatches, and it stays discharged only when all cells match~\cite{levi2026gainp}.
Driving both searchlines of a row high masks the row.
Storing 1 in both gain cells of a CAM cell blocks both read paths for any query, which realizes a ternary ``don't care''.
The \gainp{} TCAM cell thus needs six transistors~\cite{levi2026gainp}.

\begin{figure}[t]
\centering
\resizebox{\textwidth}{!}{%
\begin{circuitikz}[font=\small]
% ---------------- (a) hybrid-3T gain cell ----------------
\begin{scope}
  \draw (0,2.0) -- (0,-3.2) node[below] {WBL};
  \node[nmos, rotate=-90] (MW) at (1.0,0) {};
  \draw (0,0) -- (MW.S);
  \draw (MW.G) -- ++(0,0.5) node[above] {WWL};
  \node[below=1pt] at (MW.center) {\footnotesize M$_\mathrm{W}$};
  \draw (MW.D) -- (2.2,0) node[circ] (sn) {};
  \node[above] at (2.2,0.05) {SN};
  \node[pmos] (MS) at (3.4,0) {};
  \draw (2.2,0) -- (MS.G);
  \draw (MS.S) -- ++(0,0.25) node[vcc] {};
  \node[right] at (MS.east) {\footnotesize M$_\mathrm{S}$};
  \node[pmos] (MR) at (3.4,-2.2) {};
  \draw (MS.D) -- (MR.S);
  \draw (MR.G) -- ++(-0.35,0) node[left] {RWL};
  \node[right] at (MR.east) {\footnotesize M$_\mathrm{R}$};
  \draw (MR.D) -- ++(0,-0.35) -| (4.9,-2.6);
  \draw (4.9,2.0) -- (4.9,-3.2) node[below] {RBL};
  \draw (1.85,0) node[circ] {} to[C, l^=\footnotesize$C_\mathrm{SN}$, bipoles/length=0.8cm] (1.85,-1.1) node[ground] {};
  \node at (2.4,-5.9) {(a)};
\end{scope}
% ---------------- (b) CAM cell, column, sense and tag ----------------
\begin{scope}[xshift=11.4cm, every node/.style={font=\small}]
  \tikzset{blk/.style={draw, rounded corners=2pt, minimum height=0.85cm, align=center, fill=white}}
  \node[blk, minimum width=2.9cm] (gt) at (0,1.5) {gain cell, read port\\ $\mathrm{SN}=\beta$};
  \node[blk, minimum width=2.9cm] (gb) at (0,0.0) {gain cell, read port\\ $\mathrm{SN}=\overline{\beta}$};
  \draw[dashed, rounded corners] (-1.8,-0.75) rectangle (1.8,2.3);
  \node[anchor=south east, font=\footnotesize] at (1.8,2.3) {CAM cell};
  \draw[-{Latex}] (-3.3,1.5) node[left, align=right] {$\overline{\mathrm{SL}}$\\(top RWL)} -- (gt.west);
  \draw[-{Latex}] (-3.3,0.0) node[left, align=right] {SL\\(bottom RWL)} -- (gb.west);
  \draw (2.3,2.6) -- (2.3,-1.4);
  \draw (gt.east) -- (2.3,1.5) node[circ] {};
  \draw (gb.east) -- (2.3,0.0) node[circ] {};
  \node[right, align=left, font=\footnotesize] at (2.35,2.4) {ML\\(RBL)};
  \node[font=\footnotesize, align=center] at (0,-1.25) {other CAM cells\\of the column};
  \draw[dotted, thick] (0,-1.72) -- (0,-1.92);
  \node[blk, minimum width=1.6cm] (sa) at (2.3,-2.2) {SA};
  \draw (2.3,-1.4) -- (sa.north);
  \node[blk, minimum width=1.6cm] (or) at (2.3,-3.5) {OR};
  \node[blk, minimum width=1.6cm] (ff) at (2.3,-4.8) {tag FF};
  \draw[-{Latex}] (sa) -- (or);
  \draw[-{Latex}] (or) -- (ff);
  \draw[-{Latex}] (ff.east) -- ++(0.5,0) |- (or.east);
  \draw[-{Latex}] (ff.west) -- ++(-1.2,0) node[left, align=right, font=\footnotesize] {to write\\bitlines};
  \node[anchor=west, align=left, font=\footnotesize] at (3.4,-2.2) {match / mismatch};
  \node[anchor=west, align=left, font=\footnotesize] at (3.6,-4.4) {one-bit\\accumulator};
  \node at (1.0,-5.9) {(b)};
\end{scope}
\end{circuitikz}}
\caption{(a) Hybrid-3T gain cell. (b) \gainp{} CAM cell built from two stacked gain cells (read ports only), the matchline, the sense amplifier and the tag register/accumulator of one column. Redrawn after the description in~\cite{levi2026gainp}.}
\label{fig:gc}
\end{figure}

\paragraph{Microarchitecture}
The array consists of subarrays of $256\times64$ gain cells, i.e., $128\times64$ CAM cells (Fig.~\ref{fig:arch}(a))~\cite{levi2026gainp}.
Each subarray has the gain-cell core, row logic that drives the WWLs and RWLs, sense amplifiers and tag register/accumulators.
The tag of a column ORs the sense amplifier output with its previous value. It thus accumulates matches over consecutive compare cycles and is reset after a write~\cite{levi2026gainp}.
The tags drive the write bitlines, so a write stores the tag value into a selected row of every column.
A unified controller issues the compare and write sequences to all subarrays~\cite{levi2026gainp}.

\paragraph{Associative computing}
An AP computes by compare and write steps applied to all columns in parallel.
An arbitrary Boolean function of a few rows costs one compare per product term of its sum-of-products form, followed by one write.
For example, \gainp{} vector addition takes $5+4=9$ cycles per bit: four compares and a write for the sum, three compares and a write for the carry~\cite{levi2026gainp}.
The cost depends on the operand width, not on the number of columns.

\paragraph{Circuit figures and refresh}
The 16~nm gain cell measures 0.082~$\mu$m$^2$, and the $256\times64$ macro 1889.9~$\mu$m$^2$, an array efficiency of 75\%~\cite{levi2026gainp}.
Post-layout simulation with process variations gives a 2~ns compare cycle and 10.4~fJ per bit per compare. The evaluation phase accounts for 42\% of the delay and 45\% of the power~\cite{levi2026gainp}.
A 1000-point worst-case Monte Carlo simulation gives a retention time of 75.67~$\mu$s mean and 7.57~$\mu$s standard deviation~\cite{levi2026gainp}.
More mismatching cells per column widen the separation between match and mismatch matchline voltages~\cite{levi2026gainp}.

\begin{figure}[t]
\centering
\resizebox{0.95\textwidth}{!}{%
\begin{tikzpicture}[font=\small, >={Latex}]
\tikzset{sub/.style={draw, fill=blue!6, minimum width=1.45cm, minimum height=0.9cm, font=\scriptsize, align=center},
         blk/.style={draw, rounded corners=2pt, align=center, fill=white}}
% (a) array of subarrays and controller
\node[blk, minimum width=6.6cm, minimum height=0.7cm] (ctl) at (2.4,4.1) {controller (in this work: sequencer and model SRAM)};
\foreach \i in {0,1,2,3} \foreach \j in {0,1,2} {
  \node[sub] (s\i\j) at (\i*1.6, 2.6-\j*1.1) {subarray\\[-2pt]$128{\times}64$};
}
\draw[->] (ctl.south) -- (2.4,3.05);
\node[font=\footnotesize, align=center] at (2.4,-0.65) {each subarray: GC core, row logic (WWL, RWL drivers),\\ sense amplifiers, tag register/accumulators};
\node at (2.4,-1.5) {(a)};
% (b) logical column built from stacked subarrays
\begin{scope}[xshift=9.0cm]
  \foreach \k/\y in {0/3.4, 1/2.1, 2/0.8} {
    \node[sub, minimum width=2.2cm, minimum height=1.0cm] (c\k) at (0,\y) {subarray $\k$\\[-2pt]128 CAM rows};
  }
  \foreach \y in {3.4, 2.1, 0.8} {\draw[very thick, red!70!black] (1.1,\y) -- (1.4,\y);}
  \node[font=\footnotesize] at (0,-0.15) {\vdots};
  \draw[very thick, red!70!black] (1.4,3.4) -- (1.4,-0.6) -- (1.75,-0.6);
  \node[blk, minimum width=1.0cm] (sa) at (2.3,-0.6) {SA};
  \node[blk, minimum width=1.2cm] (tag) at (2.3,-1.7) {tag};
  \draw[->] (sa) -- (tag);
  \node[anchor=west, align=left, font=\footnotesize, text width=4.6cm] at (1.6,2.75) {read bitlines of the stacked\\subarrays, joined at the input\\of one SA during COMPARE\\(switchable)};
  \node[anchor=west, align=left, font=\footnotesize] at (3.0,-1.2) {one logical column\\$= h$ stacked subarrays};
  \node at (0.6,-2.4) {(b)};
\end{scope}
\end{tikzpicture}}
\caption{(a) \gainp{} array organization, after~\cite{levi2026gainp}. (b) Logical column used in this work. During a compare, the read bitlines of $h$ vertically stacked subarrays are joined at the input of one sense amplifier. One compare thus spans all $h$ subarrays (one column shown). The local sense amplifiers of each subarray, used for refresh, are not shown.}
\label{fig:arch}
\end{figure}

Table~\ref{tab:gainp} lists the \gainp{} figures used in this work.
The article also reports a power of 2.17~$\mu$W per bit at an operating frequency of 1~GHz, evaluated with a combination of silicon measurements and circuit simulation~\cite{levi2026gainp}.
The evaluation workbook of that study models the system at the same 1~GHz clock. This 1~ns cycle is shorter than the simulated 2~ns compare cycle.
The static power of 380~$\mu$W per Mbit is taken from the same workbook, which is not published. \hyp{We therefore treat it as an assumption.}
Section~\ref{sec:gainp7} derives the predicted 7~nm parameters from these two pairs of 16~nm figures.

\begin{table}[!htbp]
\small
\centering
\caption{\gainp{} parameters used in this work, from~\cite{levi2026gainp}. The system model clock and the static power are taken from the unpublished evaluation workbook of the same study. The static power is treated as an assumption.}
\label{tab:gainp}
\begin{tabular}{ll}
\toprule
Parameter & Value \\
\midrule
Technology & 16~nm FinFET, 0.8~V \\
Subarray & $256\times64$ gain cells \\
 & ($128\times64$ CAM cells) \\
Gain cell area & 0.082~$\mu$m$^2$ \\
Subarray (macro) area & 1889.9~$\mu$m$^2$ \\
Compare cycle, circuit simulation & 2~ns \\
Compare energy, circuit simulation & 10.4~fJ per bit \\
System model clock & 1~GHz \\
System model power & 2.17~$\mu$W per bit at 1~GHz \\
Static power (unpublished workbook) & \hyp{380~$\mu$W per Mbit} \\
Retention (worst-case Monte Carlo) & 75.67~$\mu$s mean, 7.57~$\mu$s std.\ \\
\bottomrule
\end{tabular}
\end{table}
